% Problem record op_37f7b003ec78a028. This TeX file is derived from
% database/problems_json/op_37f7b003ec78a028.json by scripts/sync-tex.mjs; edit the JSON record
% and rerun the script rather than editing this file.
\section{Asymptotic growth of the stabilizer rank of T-state tensor powers}
\label{sec:op_37f7b003ec78a028}

\paragraph{Problem.}
Does the exact stabilizer rank of the tensor powers of the single-qubit
magic state $|T\rangle$ grow polynomially in the number of copies, or is
it not polynomially bounded?  The state is
\begin{equation}
  |T\rangle := 2^{-1/2}\bigl(|0\rangle + e^{i\pi/4}|1\rangle\bigr),
  \label{eq:srtk-t-state}
\end{equation}
and, for an $n$-qubit pure state $|\psi\rangle$, its exact stabilizer rank
is
\begin{equation}
  \chi(\psi) := \min\Bigl\{r : |\psi\rangle = \sum_{j=1}^{r} c_j\,|\varphi_j\rangle
  \text{ for some } c_j \in \mathbb{C} \text{ and } n\text{-qubit stabilizer states }
  |\varphi_j\rangle\Bigr\},
  \label{eq:srtk-rank-def}
\end{equation}
where an $n$-qubit stabilizer state is the common $+1$ eigenstate of $n$
independent commuting Hermitian Pauli operators.  Which of the two alternatives holds for
$\chi\bigl(|T\rangle^{\otimes n}\bigr)$: does there exist a constant $k$
with $\chi\bigl(|T\rangle^{\otimes n}\bigr) \le n^{k}$ for all sufficiently large $n$, or is
$\chi\bigl(|T\rangle^{\otimes n}\bigr)$ not bounded by any polynomial in
$n$?  The best current bounds are
\begin{equation}
  \Omega\!\left(\frac{n^{2}}{\operatorname{polylog} n}\right) \le
  \chi\bigl(|T\rangle^{\otimes n}\bigr) \le O\bigl(2^{\alpha n}\bigr),
  \qquad \alpha = \tfrac{1}{4}\log_{2} 3 \le 0.3963,
  \label{eq:srtk-bounds}
\end{equation}
which leave both alternatives open.  A sharper sub-question is whether the
exponential rate $\gamma = \lim_{n\to\infty} \frac{1}{n}\log_{2}
\chi\bigl(|T\rangle^{\otimes n}\bigr)$, which exists by sub-multiplicativity
of the rank under tensor product and Fekete's lemma, is strictly positive,
since Eq.~\eqref{eq:srtk-bounds} only pins $\gamma$ to the interval
$[0, 0.3963]$.  The small-copy frontier of the quantity in
Eq.~\eqref{eq:srtk-rank-def} on copies of the state in
Eq.~\eqref{eq:srtk-t-state} is also open: the exact values $\chi\bigl(|T\rangle^{\otimes 2}\bigr) = 2$
and $\chi\bigl(|T\rangle^{\otimes 3}\bigr) = 3$ are proved, while at four
through eight copies only upper bounds are known,
$\chi\bigl(|T\rangle^{\otimes 4}\bigr) \le 4$,
$\chi\bigl(|T\rangle^{\otimes 5}\bigr) \le 6$,
$\chi\bigl(|T\rangle^{\otimes 6}\bigr) \le 6$, and
$\chi\bigl(|T\rangle^{\otimes 7}\bigr) \le 12$, and $\chi\bigl(|T\rangle^{\otimes 8}\bigr) \le 12$ as tabulated by Qassim, Pashayan, and Gosset, the four-copy bound having
since been certified as the exact value $\chi\bigl(|T\rangle^{\otimes
4}\bigr) = 4$ by an exhaustive computer search; determine the exact values
at $t = 5$, $6$, $7$, and $8$.

\subsection*{Status}

Unsolved

\subsection*{Source}

Bravyi, Smith, and Smolin introduced the quantity
$\chi\bigl(|T\rangle^{\otimes n}\bigr)$ and its use for classical
simulation, explicitly posed its asymptotic scaling, observed that
super-polynomial growth follows assuming universal quantum computation
cannot be simulated classically in polynomial time, proved only an
$\Omega(\sqrt{n}\,)$ lower bound while stating that they were unable to
derive an unconditional super-polynomial one, and conjectured
$\chi\bigl(|T\rangle^{\otimes n}\bigr) \ge 2^{\Omega(n)}$
\sourcecite{ref:srtk-bss16}{BSS16}; Bravyi, Browne, Calpin, Campbell,
Gosset, and Howard sharpened the question to its polynomial-or-not form,
stating that polynomial scaling of the exact or approximate stabilizer
rank would entail complexity-theoretic heresies such as
$\mathsf{BQP} = \mathsf{BPP}$ or $\mathsf{P} = \mathsf{NP}$, and that no
techniques are known for proving unconditional super-polynomial lower
bounds \sourcecite{ref:srtk-bbc19}{BBC+19}.

\subsection*{Progress}

\begin{itemize}
  \item Bravyi, Smith, and Smolin introduced low-rank stabilizer decompositions
  of magic-state tensor powers for classical simulation and proved an
  $\Omega(\sqrt{n}\,)$ lower bound on the exact stabilizer rank of tensor
  powers of a single-qubit magic state, finding
  $\chi\bigl(|T\rangle^{\otimes 6}\bigr) = 7$ numerically, a value
  conjectured optimal \sourcecite{ref:srtk-bss16}{BSS16}; the sampling algorithm of Bravyi and
  Gosset for circuits dominated by Clifford gates had used approximate
  stabilizer decompositions \sourcecite{ref:srtk-bg16}{BG16}.

  \item Bravyi, Browne, Calpin, Campbell, Gosset, and Howard developed the theory
  of exact and approximate stabilizer rank; the proved exact values are
  $\chi\bigl(|T\rangle^{\otimes 2}\bigr) = 2$ and
  $\chi\bigl(|T\rangle^{\otimes 3}\bigr) = 3$, and past three copies only
  upper bounds are known, namely $\chi\bigl(|T\rangle^{\otimes 4}\bigr)
  \le 4$, $\chi\bigl(|T\rangle^{\otimes 5}\bigr) \le 6$, and
  $\chi\bigl(|T\rangle^{\otimes 7}\bigr) \le 12$ as tabulated by Qassim,
  Pashayan, and Gosset, and no finite set of copy counts can settle the
  asymptotic growth in either direction
  \sourcecite{ref:srtk-bbc19}{BBC+19},
  \sourcecite{ref:srtk-qpg21}{QPG21}.

  \item The upper bound in Eq.~\eqref{eq:srtk-bounds} is due to Qassim, Pashayan,
  and Gosset: entangled six-qubit magic cat states of stabilizer rank
  exactly $3$ give $\chi\bigl(|T\rangle^{\otimes n}\bigr) \le
  O\bigl(2^{\alpha n}\bigr)$, with the exponent attained only in the limit
  of an infinitely long chain of contracted cat states, improving the previous
  $\chi\bigl(|T\rangle^{\otimes 12}\bigr) \le 47$, that is $\alpha \le
  0.463$, of Kocia \sourcecite{ref:srtk-koc20}{Koc20}, and disproving the
  conjectured value $7$ at six copies by showing
  $\chi\bigl(|T\rangle^{\otimes 6}\bigr) \le 6$
  \sourcecite{ref:srtk-qpg21}{QPG21}.

  \item Peleg, Shpilka, and Volk proved an $\Omega(n)$ lower bound on the exact
  stabilizer rank of tensor powers of single-qubit magic states, improving
  the $\Omega(\sqrt{n}\,)$ of Bravyi, Smith, and Smolin, as well as
  $\Omega(\sqrt{n}/\log n)$ for $\delta$-close approximate rank, the
  first nontrivial approximate bound; these direct exact-rank techniques
  stop at linear \sourcecite{ref:srtk-psv22}{PSV22}.

  \item Two further exact-rank techniques reach only linear lower bounds:
  higher-order Fourier analysis \sourcecite{ref:srtk-lab22}{Lab22}, and the
  new techniques of Lovitz and Steffan
  \sourcecite{ref:srtk-ls22}{LS22}.

  \item Mehraban and Tahmasbi proved $\chi_{\delta}\bigl(|T\rangle^{\otimes
  m}\bigr) = \Omega\bigl((1-\delta^{2})^{2} m^{2}/\operatorname{polylog}
  m\bigr)$ for every $0 < \delta < 1$, where $\chi_{\delta}(\psi)$ is the
  least stabilizer rank of a pure state within $2$-norm distance $\delta$
  of $|\psi\rangle$; fixing any constant $\delta$ and using the
  monotonicity $\chi \ge \chi_{\delta}$ yields the lower bound in
  Eq.~\eqref{eq:srtk-bounds}, the only super-linear technique known.  They
  further show that polynomial exact rank would imply
  $\mathsf{P}^{\#\mathsf{P}} \subseteq \mathsf{P}/\mathsf{poly}$
  (Theorem~1.6), give a conditional super-polynomial approximate-rank bound
  under average-case $\#\mathsf{P}$-hardness (Theorem~1.11), and prove that
  a quadratic approximate bound alone cannot yield super-quadratic exact
  bounds in general (Lemma~A.5) \sourcecite{ref:srtk-mt24}{MT24}.

  \item Kalra and Sinha proved the first quantitative lower bound on the
  stabilizer fidelity as a function of stabilizer rank, via Barnes-Wall
  lattices and magic monotones, but for stabilizer rank it reproduces only
  the existing lower bound, linear up to a logarithmic factor, on the
  approximate rank of the tensor powers of the Hadamard magic state
  $|H\rangle$, together with an elementary proof of the Lovitz-Steffan
  density of maximal-stabilizer-rank product states; neither side of
  Eq.~\eqref{eq:srtk-bounds} for $|T\rangle^{\otimes n}$ moves
  \sourcecite{ref:srtk-ks25}{KS25}.

  \item Labib and Russo gave explicit stabilizer decompositions for the three
  previously unbounded qutrit magic-state orbits, with per-copy exponents
  $0.316$ and $0.421$, and the first $\Omega(m/\log m)$ lower bounds for
  the qutrit Hadamard-eigenstate and Norrell orbits; for qubits, a
  closed-form four-copy decomposition of stabilizer rank exactly $3$ for
  the face-center orbit (T-type in their naming, which swaps T and H
  relative to this record: their H-type orbit contains the $|T\rangle$ of
  Eq.~\eqref{eq:srtk-t-state}) merely matches the exponent $\alpha$ of
  Eq.~\eqref{eq:srtk-bounds}, and an exhaustive search there rigorously
  confirms $\chi\bigl(|T\rangle^{\otimes 4}\bigr) = 4$, so neither side
  of the qubit question moves \sourcecite{ref:srtk-lr26}{LR26}.
\end{itemize}

\subsection*{References}

\begin{enumerate}[label={},leftmargin=4.8em,labelsep=0.5em]
  \footnotesize
  \item[\textup{[BSS16]}]\label{ref:srtk-bss16}
  S. Bravyi, G. Smith, and J. A. Smolin, ``Trading classical and quantum
  computational resources,'' \emph{Physical Review X} \textbf{6}, 021043
  (2016).
  \href{https://doi.org/10.1103/PhysRevX.6.021043}{doi:10.1103/PhysRevX.6.021043};
  \href{https://arxiv.org/abs/1506.01396}{arXiv:1506.01396}.

  \item[\textup{[BG16]}]\label{ref:srtk-bg16}
  S. Bravyi and D. Gosset, ``Improved classical simulation of quantum
  circuits dominated by Clifford gates,'' \emph{Physical Review Letters}
  \textbf{116}, 250501 (2016).
  \href{https://doi.org/10.1103/PhysRevLett.116.250501}{doi:10.1103/PhysRevLett.116.250501}.

  \item[\textup{[BBC+19]}]\label{ref:srtk-bbc19}
  S. Bravyi, D. Browne, P. Calpin, E. Campbell, D. Gosset, and M. Howard,
  ``Simulation of quantum circuits by low-rank stabilizer decompositions,''
  \emph{Quantum} \textbf{3}, 181 (2019).
  \href{https://doi.org/10.22331/q-2019-09-02-181}{doi:10.22331/q-2019-09-02-181};
  \href{https://arxiv.org/abs/1808.00128}{arXiv:1808.00128}.

  \item[\textup{[Koc20]}]\label{ref:srtk-koc20}
  L. Kocia, ``Improved strong simulation of universal quantum circuits,''
  \href{https://arxiv.org/abs/2012.11739}{arXiv:2012.11739} (2020).

  \item[\textup{[QPG21]}]\label{ref:srtk-qpg21}
  H. Qassim, H. Pashayan, and D. Gosset, ``Improved upper bounds on the
  stabilizer rank of magic states,'' \emph{Quantum} \textbf{5}, 606
  (2021).
  \href{https://doi.org/10.22331/q-2021-12-20-606}{doi:10.22331/q-2021-12-20-606};
  \href{https://arxiv.org/abs/2106.07740}{arXiv:2106.07740}.

  \item[\textup{[PSV22]}]\label{ref:srtk-psv22}
  S. Peleg, A. Shpilka, and B. L. Volk, ``Lower Bounds on Stabilizer
  Rank,'' \emph{Quantum} \textbf{6}, 652 (2022).
  \href{https://doi.org/10.22331/q-2022-02-15-652}{doi:10.22331/q-2022-02-15-652};
  \href{https://arxiv.org/abs/2106.03214}{arXiv:2106.03214}.

  \item[\textup{[Lab22]}]\label{ref:srtk-lab22}
  F. Labib, ``Stabilizer rank and higher-order Fourier analysis,''
  \emph{Quantum} \textbf{6}, 645 (2022).
  \href{https://doi.org/10.22331/q-2022-02-09-645}{doi:10.22331/q-2022-02-09-645}.

  \item[\textup{[LS22]}]\label{ref:srtk-ls22}
  B. Lovitz and V. Steffan, ``New techniques for bounding stabilizer
  rank,'' \emph{Quantum} \textbf{6}, 692 (2022).
  \href{https://doi.org/10.22331/q-2022-04-20-692}{doi:10.22331/q-2022-04-20-692}.

  \item[\textup{[MT24]}]\label{ref:srtk-mt24}
  S. Mehraban and M. Tahmasbi, ``Quadratic Lower bounds on the Approximate
  Stabilizer Rank: A Probabilistic Approach,'' in \emph{Proceedings of
  STOC 2024} (2024).
  \href{https://doi.org/10.1145/3618260.3649733}{doi:10.1145/3618260.3649733};
  \href{https://arxiv.org/abs/2305.10277}{arXiv:2305.10277}.

  \item[\textup{[KS25]}]\label{ref:srtk-ks25}
  A. R. Kalra and P. Sinha, ``Stabilizer Ranks, Barnes Wall Lattices and
  Magic Monotones,''
  \href{https://arxiv.org/abs/2503.04101}{arXiv:2503.04101} (2025).

  \item[\textup{[LR26]}]\label{ref:srtk-lr26}
  F. Labib and V. Russo, ``Stabilizer rank bounds for magic-state
  orbits,'' \href{https://arxiv.org/abs/2605.28586}{arXiv:2605.28586}
  (2026).
\end{enumerate}

\subsection*{Comment}

What remains unresolved is the alternative itself: neither an $n^{O(1)}$
upper bound nor a super-polynomial lower bound on
$\chi\bigl(|T\rangle^{\otimes n}\bigr)$ is known, and the two sides of
Eq.~\eqref{eq:srtk-bounds} have stood since 2021 and 2024 respectively.
A polynomial upper bound established merely as a rank statement would
not by itself give a uniform simulation: a bound on $\chi$ does not
construct a stabilizer decomposition of $|T\rangle^{\otimes m}$ or compute
its coefficients, and by gadgetizing every T gate through state injection
into a Clifford circuit acting on $|T\rangle^{\otimes m}$ with the
decomposition supplied as polynomial advice, Theorem~1.6 of
\sourcecite{ref:srtk-mt24}{MT24} draws from such a bound only the
nonuniform consequence
$\mathsf{P}^{\#\mathsf{P}} \subseteq \mathsf{P}/\mathsf{poly}$, so it is
believed false.  A polynomial-time classical strong-simulation algorithm
for all Clifford+T circuits with polynomially many T gates would need the
additional hypothesis that some polynomial-length decomposition of
$|T\rangle^{\otimes m}$ can be constructed, with coefficients computable
to the required precision, in polynomial time, as the explicit cat-state
constructions behind Eq.~\eqref{eq:srtk-bounds} can be.  Settling the super-polynomial alternative requires a lower bound that
rules out $\chi\bigl(|T\rangle^{\otimes n}\bigr) \le n^{k}$ for every
constant $k$; a bound of one fixed exponent, such as $n^{2+\varepsilon}$
for some $\varepsilon > 0$, would improve the current
$\Omega\bigl(n^{2}/\operatorname{polylog} n\bigr)$ but is still polynomial
and would leave the alternative open.  A genuinely super-polynomial lower
bound would be the first unconditional bound beyond the
$\Omega\bigl(n^{2}/\operatorname{polylog} n\bigr)$ level, since the
direct methods of Peleg, Shpilka, and Volk, Labib, and Lovitz and Steffan
stop at linear and a quadratic approximate bound alone cannot give
super-quadratic exact bounds in general.  Constant-factor improvements ---
a smaller $\alpha$ in Eq.~\eqref{eq:srtk-bounds}, sharpening the lower
bound to $\Omega(n^{2})$, or new small-$n$ exact values --- do not settle
the question.  The zoo's record ``Universal purification with classically
simulable operations'' shares the boundary between stabilizer and magic
resources but asks a dynamical resource-conversion question about
purification fidelity, not a state-decomposition question.

\subsection*{Field}

Quantum algorithm

\subsection*{Topic}

Quantum magic; Quantum circuit complexity; Computational complexity and computability

\subsection*{ID}

\texttt{op\_37f7b003ec78a028}
